% MT060 / explanation — canonical editable content.
记软间隔SVM的目标值为\(J_C(w,b)\)：



\[
J_C(w,b)=\frac12\|w\|^2+C\sum_i\max(0,1-y_if(x_i)).
\]


第一项偏好较小权重、较宽间隔；第二项对分错或离边界太近的点收费。\(C\)乘在违例损失前，所以越大，模型越不愿意接受个别点的违例。


\textbf{(a) 大\(C\)的草图。} 先看两团主数据，再看偏离主团的红点。若二次特征空间能把它们分开，大\(C\)会更努力照顾这些点，边界因而弯曲。原解的左图试图把那两个红点也留在红类一侧。不能泛化成“任何核、任何数据、\(C\)大都必然零错误”。


\textbf{(b) 先回答原题的严格极限。} 这里各条样本同等计入损失，偏置\(b\)不受正则惩罚，核函数固定。原图有10个红点、7个绿点；令红类标签为\(-1\)、绿类为\(+1\)。分数写成\(f(x)=w^T\phi(x)+b\)，其中\(\phi(x)\)是二次核对应的特征，预测类别由分数正负决定。


先尝试一个很简单的模型：\(w=0,b=-1\)，所有输入都得到分数\(-1\)，全部判红。红点的\(yf(x)=1\)，hinge损失为0；绿点的\(yf(x)=-1\)，每个损失为\(1-(-1)=2\)。7个绿点的总损失为14，目标值就是\(14C\)。


最优解不会比这个候选解更差。目标中的hinge项又不可能为负，所以最优权重\(w_C\)满足



\[
\frac12\|w_C\|^2\le J_C(w_C,b_C)\le14C,
\qquad \|w_C\|^2\le28C.
\]


因此\(C\to0^+\)时，\(w_C\)也趋于0，固定输入的分数越来越接近同一个偏置\(b\)。偏置该取多少，要继续比较两类的hinge损失。忽略正在趋于0的内积分数后，总损失为



\[
H(b)=10\max(0,1+b)+7\max(0,1-b).
\]


当\(-1\le b\le1\)时，\(H(b)=17+3b\)，\(b\)越小损失越低，最低在\(b=-1\)。低于\(-1\)后，红点已不再获益，绿点的损失却继续上升；高于1也只会让红点损失增大。因此最优偏置趋于\(-1\)，图示区域最终全部判为红类，不再有分开两团点的边界。画严格极限时，应把整个区域标为红类。


\textbf{随卷右图对应的是较小的正\(C\)。} 在还未趋于零的某个正\(C\)下，模型可以接受少数点违例，在主团之间保留分界。这幅原解草图仍可帮助理解取舍，但不能当作严格极限。原题用红圆点和绿星形区分类别；答案草图把两类都画成圆点，建议保留彩色。


\textbf{(c) 测试表现取决于离群点的性质。} 若它们是标签噪声，取适中的正\(C\)、允许这些点分错，可能比追着它们拟合更好；若它们代表真实但少见的规律，忽略它们也会犯错。另一方面，严格趋零会把全部绿类丢掉，不能仅凭“正则更强”就说它最好。应在验证数据上选\(C\)。


选读：对有限训练点，令\(s_i=w_C^T\phi(x_i)\)，则所有\(s_i\to0\)。固定\(w_C\)时，红点的hinge折点为\(b=-1-s_i\)。\(b\)低于全部红点折点时，总损失随\(b\)增大而下降；越过全部红点折点、还没碰到绿点折点时，总斜率为\(10-7=3>0\)。因此最优\(b_C\)夹在这些趋于\(-1\)的红点折点之间，确实有\(b_C\to-1\)。若直接设\(C=0\)，整个数据项消失，\(b\)可以任取；这与从正数一侧趋近0不同。


\StudyImage{CS5489/exam-review/Midterm/assets/answer-mt060-pair.png}


\emph{左：原答案的大C草图；右：较小正C草图，不能代表严格\(C\to0^+\)时全判红类的极限。}


\textbf{原题差异 / Source note：} 原解把较小正C的分隔曲线用于C趋零；按本题标准目标，严格极限全判多数红类。 / The supplied small-positive-C sketch is not the strict limit, which predicts the majority red class under the stated standard objective.


\textbf{出处：} 2023B期中 Q9，原卷第6页（10分）


\textbf{答案依据：} 2023B期中 Q9 随卷／配套参考答案；条件与解释经复核。
